\section{Stationary estimates for derivative growth}\label{sec:stationary}

The switching construction requires control of derivatives on a long finite
orbit segment.  The accuracy may depend on the length of the segment, but
the squared errors must be summable over dyadic lengths.  We first obtain
this summability for an arbitrary bounded stationary process.  Applying
the result to a triangular representation of the derivative then gives
large sets of orbit segments with the required control.

\subsection{Deviation from a dyadic chord}

Let $(Y_t)_{t\in\Z}$ be a real stationary process with
$\E Y_0^2<\infty$.  Stationarity means that translating all indices in
any finite list does not change its joint distribution.  For $h\geq 0$
define the block average and normalized maximal chord deviation by
\[
 A_h=2^{-h}\sum_{t=0}^{2^h-1}Y_t,
 \qquad
 Z_h=2^{-h}\max_{0\leq u\leq 2^h}
 \left|\sum_{t=0}^{u-1}Y_t-uA_h\right|.
\]
The maximum is over integer $u$, and an empty sum is zero.
In particular $Z_0=0$.  The chord here uses the average on the
particular block under consideration; it is not a line whose slope is
the mean of the stationary process.

\begin{lemma}[Dyadic chord estimate]\label{lem:dyadic}
For the process above, put
\[
 v_l=\E A_{l-1}^2-\E A_l^2,\qquad l\geq 1.
\]
Then $v_l\geq 0$, $\sum_{l\geq 1}v_l\leq\E Y_0^2$, and
\begin{equation}\label{eq:dyadic-convolution}
 \|Z_h\|_{L^2}
 \leq \frac12\sum_{a=0}^{h-1}2^{-a/2}\sqrt{v_{h-a}}
 \qquad(h\geq 1).
\end{equation}
Consequently
\begin{equation}\label{eq:dyadic-square-sum}
 \sum_{h\geq 1}\|Z_h\|_{L^2}^2
 \leq C\,\E Y_0^2
\end{equation}
for an absolute constant $C$.  The same estimates hold for a block
beginning at any deterministic integer time.
\end{lemma}

\begin{proof}
Consider a block of length $2^l$, and let $U$ and $V$ be the averages on
its left and right halves.  Its own average is $W=(U+V)/2$.  The identity
\[
 \frac{(U-W)^2+(V-W)^2}{2}
   =\frac{U^2+V^2}{2}-W^2
\]
and stationarity give
\[
 v_l
 =\E\frac{(U-W)^2+(V-W)^2}{2}
 =\E(U-W)^2.
\]
The last equality also follows pointwise from $V-W=-(U-W)$.
Thus $v_l\geq0$, and telescoping the differences of second moments proves
the asserted bound on their sum.  No independence of the two halves is
used.

We next compare the cumulative-sum polygon on $[0,2^h]$ with its chord.
For $a=0,\ldots,h$, let $F_a$ be the piecewise-linear interpolant that
agrees with the cumulative sum at the multiples of $2^{h-a}$.
Then $F_0$ is the chord and $F_h$ agrees with every integer-time
cumulative sum.  Passing from $F_a$ to $F_{a+1}$ inserts the midpoint
of each of $2^a$ intervals of length $2^{h-a}$.
If $U_j,W_j$ are respectively the left-child and parent averages on
the $j$th such interval, the displacement at its midpoint is
$2^{h-a-1}(U_j-W_j)$.
The difference of the interpolants is linear between the old and new
vertices, so
\[
 2^{-h}\|F_{a+1}-F_a\|_\infty
 \leq 2^{-a-1}\max_{1\leq j\leq 2^a}|U_j-W_j|.
\]
Bounding a squared maximum by the sum of squares and using stationarity
of each interval gives
\[
 \left\|2^{-h}\|F_{a+1}-F_a\|_\infty\right\|_{L^2}
 \leq 2^{-a-1}\bigl(2^a v_{h-a}\bigr)^{1/2}.
\]
The triangle inequality, first for the polygonal sup norm and then in
$L^2$, proves \eqref{eq:dyadic-convolution}.

Extend the sequence $(\sqrt{v_l})_{l\geq1}$ by zero to nonpositive
indices.  Its squared $\ell^2$ norm is at most $\E Y_0^2$.
The right side of \eqref{eq:dyadic-convolution} is its convolution
with the summable sequence $(2^{-a/2-1})_{a\geq0}$.
The triangle inequality for the $\ell^2$ norms of its translates
therefore proves \eqref{eq:dyadic-square-sum}.
Finally, stationarity identifies the distribution of each translated
block with that of the block just considered.
\end{proof}

\subsection{A stationary triangular derivative representation}

Fix a compact smooth manifold $M$ of dimension $d\geq1$, a smooth
Riemannian metric, a smooth diffeomorphism $S:M\to M$, and an
$S$-invariant probability measure $\nu$.  In the switching argument
$\nu$ will have a smooth positive density.  That extra assumption is
not needed in this section.  The constructions also apply to a compact
manifold with smooth boundary.

The compact frame extension and triangular cocycle have a classical
precedent in Oseledets's proof of the multiplicative ergodic theorem
\cite[\S4, pp.~208--209, equations~(6)--(10)]{Oseledets1968}.
We use the full orthogonal frame bundle and a positive-diagonal upper
triangular convention. The finite-time estimates below follow from the
stationary chord lemma; we do not invoke a limiting Lyapunov splitting
or a convergence rate in the multiplicative ergodic theorem.

The orthonormal frame bundle $\mathcal F M$ consists of the linear
isometries $e:\R^d\to T_xM$, with bundle projection $\pi(e)=x$.
For each frame there is a unique positive-diagonal QR factorization
\begin{equation}\label{eq:frame-qr}
 D S_x\circ e=\widetilde S(e)\,R(e),
\end{equation}
where $\widetilde S(e)$ is an orthonormal frame at $Sx$ and $R(e)$
is upper triangular with positive diagonal.
This is ordinary Gram--Schmidt in the indicated order, with the length
of each new perpendicular component chosen positive.
It depends continuously, indeed smoothly, on the frame and the base
point.

The map $\widetilde S$ is invertible.  If
$DS_x e=e'R$, then
$(DS_x)^{-1}e'=D(S^{-1})_{Sx}e'=eR^{-1}$, where the inverse
linear map has domain $T_{Sx}M$ and codomain $T_xM$.
The matrix $R^{-1}$ is again upper triangular with
positive diagonal.  Taking positive-diagonal QR after applying the
inverse derivative therefore recovers $e$.  This also proves
continuity of the inverse.  The construction uses the full orthonormal
frame bundle and does not require a global frame or an orientation.

There is a $\widetilde S$-invariant probability $\widetilde\nu$ whose
projection is $\nu$.  To see this, first put the normalized
orthogonal-group Haar measure on every fiber over $\nu$.
Orthogonal changes of local trivialization preserve this fiber
measure, so the resulting probability is defined without choosing a
global section.  Each of its iterates under $\widetilde S$ projects
to $\nu$.  Cesaro averages have a weakly convergent subsequence,
because $\mathcal F M$ is compact.  The difference between the
pushforward of an average and the average itself tends to zero
against every continuous function, so its limit is invariant.
Continuity of $\pi$ preserves the marginal $\nu$ in the limit.

Compactness and invertibility provide a constant $A_*>1$ such that
\begin{equation}\label{eq:qr-uniform-bounds}
 \|R(e)\|,\ \|R(e)^{-1}\|\leq A_*
 \qquad(e\in\mathcal F M).
\end{equation}
In particular $A_*^{-1}\leq R(e)_{jj}\leq A_*$.
The functions
\[
 g_j(e)=\ln R(e)_{jj},\qquad 1\leq j\leq d,
\]
are continuous and bounded, and
$(g_j(\widetilde S^t e))_{t\in\Z}$ is stationary under
$\widetilde\nu$.

\begin{proposition}[Good finite orbit segments]\label{prop:good-frames}
For each $\eta>0$ there is a positive sequence $(b_k)_{k\geq1}$ with
\begin{equation}\label{eq:b-summability}
 b_k\geq 2^{-k/10},
 \qquad
 \sum_{k\geq1}b_k^2<\infty,
\end{equation}
and compact sets $\widetilde{\mathcal G}_k\subset\mathcal F M$ and
$\mathcal G_k=\pi(\widetilde{\mathcal G}_k)\subset M$ such that
\[
 \widetilde\nu(\widetilde{\mathcal G}_k)\geq1-\eta/100,
 \qquad \nu(\mathcal G_k)\geq1-\eta/100.
\]
For $n=2^k$, $b=b_k$, and $e\in\widetilde{\mathcal G}_k$, define
\[
 A_t=R(\widetilde S^t e)\quad(-n\leq t<n),
 \qquad
 \lambda_j=\frac1{2n}\sum_{t=-n}^{n-1}\ln(A_t)_{jj}.
\]
Then for every $j$ and every integer $-n\leq u\leq n$,
\begin{equation}\label{eq:good-chord}
 \left|\sum_{t=-n}^{u-1}\ln(A_t)_{jj}
          -(u+n)\lambda_j\right|\leq bn.
\end{equation}
Consequently, on every subinterval $-n\leq u\leq v\leq n$,
\begin{equation}\label{eq:good-subinterval}
 \left|\sum_{t=u}^{v-1}\ln(A_t)_{jj}
          -(v-u)\lambda_j\right|\leq2bn.
\end{equation}
\end{proposition}

\begin{proof}
Apply \Cref{lem:dyadic} to each $g_j$.
Let $c_{j,h}$ be the right side of
\eqref{eq:dyadic-convolution} for that process, and put
$\alpha_h^2=\sum_{j=1}^d c_{j,h}^2$.
The lemma gives $\sum_h\alpha_h^2<\infty$.
For example, take
\[
 b_k=40\eta^{-1/2}\alpha_{k+1}+2^{-k/10}.
\]
This sequence has the properties in \eqref{eq:b-summability}.
The interval $[-n,n)$ has length $2^{k+1}$.
Stationarity and the squared Markov inequality give
\[
 \widetilde\nu\!\left(
 \max_{\substack{1\leq j\leq d\\-n\leq u\leq n}}
 \left|\sum_{t=-n}^{u-1}g_j(\widetilde S^t e)
                  -(u+n)\lambda_j\right|>b_kn
 \right)
 \leq \frac{4\alpha_{k+1}^2}{b_k^2}
 \leq\frac{\eta}{400}.
\]
Define $\widetilde{\mathcal G}_k$ by the complementary non-strict
inequalities.  These are finitely many continuous inequalities on a
compact bundle, so this set and its projection are compact.
Moreover
\[
 \pi^{-1}(M\setminus\mathcal G_k)
   \subset\mathcal F M\setminus\widetilde{\mathcal G}_k.
\]
The pushforward identity $\pi_*\widetilde\nu=\nu$ gives the asserted
base-measure estimate.  No measurable choice of a frame over each
point of $\mathcal G_k$ is needed.
Finally, subtracting the chord error at $u$ from the chord error at $v$
gives \eqref{eq:good-subinterval}, with the same rates $\lambda_j$.
\end{proof}

We call a frame in $\widetilde{\mathcal G}_k$ good at scale $k$,
and a base point in $\mathcal G_k$ good at that scale.
Every good base point has a good witness frame.  The estimates below
hold for every such witness; witnesses at different points or scales
need not be chosen coherently.  We do not require an individual point
to be good at all sufficiently large scales.

The stationary argument has now supplied the summable quantity that
will control the cost of switching.  Fix $\eta$ and its sequence
$(b_k)$, and use the following parameters at scale $n=2^k$:
\begin{equation}\label{eq:smooth-scales}
 \begin{aligned}
 b&=b_k,& B&=b^{1/8},& g&=b^{1/2},\\
 L&=\lceil b^{1/4}n\rceil,&
 h_0&=e^{-Bn},&r&=e^{-2Bn},&\delta&=e^{-3Bn/2}.
 \end{aligned}
\end{equation}
The parameter $b$ controls the diagonal chord error $bn$;
$g$ will be the separation margin for the finite-interval rates,
$L$ is the short comparison time, and $B$ sets the spatial scales
and the upper bound $CBn$ for the logarithm of the number of
partition cells. Here $h_0$ and $r$ will be spatial lengths; $\delta$ will be a
neighborhood radius for an orbit construction.
Since $b_k\to0$ and $bn\geq n^{9/10}$,
\begin{equation}\label{eq:smooth-hierarchy}
 \ln n=o(bn),\quad bn=o(gL),\quad gL=o(gn),\quad
 gn=o(L),\quad L=o(Bn).
\end{equation}
The ceiling in $L$ has no effect on these comparisons.
We also have
\[
 \sum_k B_k^{16}=\sum_k b_k^2<\infty,
 \qquad Bn\geq n^{79/80}.
\]
All subsequent asymptotic estimates are for $k\to\infty$.
Constants may depend on the fixed smooth data, dimension, and $\eta$,
but not on the point, its good witness, or the scale.
